% Stage 8 restored fragment — 2.0.2 Maxwell equations in operator form (archive lines 550-916)
\subsection{Maxwell equations in operator form} \label{sec:ch202}

A standard-gain horn on a compact-range turntable fixes the laboratory scene for `Maxwell equations in operator form'': frequency, probe standoff, and wattmeter closure on $\Gamma_R$ precede the formal displays below \cite{hansen1988}. Maxwell's retarded export to $\Gamma_R$ remains the governing constraint \cite{maxwell1873}.
Subsection~\ref{sec:ch201} named the solution map $\mathcal{S}_\omega$ and the outgoing
selector $\Pi_{\mathrm{rad}}$; the display below identifies the Maxwell operator whose
inversion, under declared boundary completion, produces that map \cite{stratton1941}. The governing question is which field pair $(\E,\H)$ is
consistent with impressed currents $\Jimp$ and with the termination class $\mathcal{B}_\Gamma$
on $\Gamma=\partial\Omega$---the same question every full-wave solver answers before any
pattern is extracted \cite{harrington2001}. Harmonic Maxwell theory on $\Omega$ is packaged
as a first-order residual system on a six-component state,
augmented by boundary traces through the operator
Eq.~\eqref{eq:ch1.boundary-data-operator} imported from
Section~\ref{sec:ch01.1.2}. Local curl and divergence balances remain those of
Eq.~\eqref{eq:ch1.maxwell-operator-block}; the Maxwell export sequence developed here does not
alter those local balances \cite{stratton1941}.
First-order Maxwell theory couples electric and magnetic fields through circulation;
packaging both components in a single state vector is the standard device that permits
resolvent notation and variational formulations without abandoning the physical coupling
\cite{stratton1941}. \(\mathbf{u}_\omega\) is
not an arbitrary six-vector: each admissible state satisfies Gauss constraints and carries
tangential traces compatible with the boundary class declared with the impressed data tuple
\(\mathfrak{D}\). From an antenna engineer's perspective, \(\mathbf{u}_\omega\) is what a
full-wave solver stores at each mesh node---the coupled \((\E,\H)\) phasor pair at fixed
\(\omega\)---before any post-processing extracts patterns or port quantities
\cite{harrington2001}. Collect phasor fields into the Maxwell state

\begin{equation}
\label{eq:ch2.maxwell-state-vector}
\mathbf{u}_\omega=
\begin{bmatrix}\E\\ \H\end{bmatrix}
\in\mathcal{F}_\Omega,
\end{equation}
where \(\mathbf{u}_\omega\) denotes the six-component Maxwell state at fixed
\(\omega\), with \(\E\) and \(\H\) carrying the SI units noted above and lying in
the admissible field space \(\mathcal{F}_\Omega\) declared in Subsection~\ref{sec:ch201}.
The pair \((\E,\H)\) is constrained jointly: divergence
rows of Eq.~\eqref{eq:ch1.maxwell-operator-block} must hold, and tangential traces on
\(\Gamma=\partial\Omega\) must lie in the domain of the boundary class declared with
\(\mathfrak{D}\). Treating \(\mathbf{u}_\omega\) as a single unknown is what permits
resolvent notation, variational forms, and the spectral compression developed in
Section~\ref{sec:ch26} \cite{harrington2001}.
Every boundary-value problem begins by distinguishing prescribed quantities from unknown
quantities. In Maxwell theory the impressed currents, charges, and boundary data constitute
the prescribed excitation, whereas the electromagnetic field remains the unknown state to
be determined \cite{maxwell1873}. The operator formulation merely records this
distinction in a compact mathematical form suitable for analysis \cite{harrington2001}.
Impressed data enter through a forcing vector aligned with the Maxwell residual rows.
The forcing packages charge and current drive on the rows where impressed sources
appear, leaving divergence constraints to be satisfied by the field unknown.
For
the tuple Eq.~\eqref{eq:ch1.data-specification-tuple},
\begin{equation}
\label{eq:ch2.maxwell-forcing-vector}
\mathbf{s}_\omega(\mathfrak{D})=
\begin{bmatrix}\mathbf{0}\\ -\Jimp\\ -\rho\\ \mathbf{0}\end{bmatrix},
\end{equation}
In Eq.~\eqref{eq:ch2.maxwell-forcing-vector},
\(\mathbf{s}_\omega(\mathfrak{D})\) is the Maxwell forcing six-vector (SI units per row)
assembled from the data tuple \(\mathfrak{D}\);
\(\Jimp\) enters the Amp\`ere row (A\,m\(^{-2}\)) and \(\rho\) the Gauss row
(C\,m\(^{-3}\)); the remaining rows are zero for the impressed-current class treated here;
row order matches Eq.~\eqref{eq:ch1.maxwell-operator-block}. Equation
\eqref{eq:ch2.maxwell-forcing-vector} records drive on the Amp\`ere and Gauss rows
only for the impressed-current class treated here; sheet and magnetic excitations
extend the vector through Eq.~\eqref{eq:ch1.interface-conditions} when trace spaces
are enriched in Section~\ref{sec:ch211} \cite{stratton1941}. The forcing is \emph{data}, not a
field variable: it specifies how \(\mathcal{S}_\omega\) is sourced while
\(\mathbf{u}_\omega\) remains the unknown. Operator language separates \emph{states}
\((\E,\H)\) from \emph{data} \((\Jimp,\rho,\mathfrak{D})\): the feed specifies what is
driven, the solver returns what Maxwell allows, and radiation operators report what is
exported \cite{maxwell1873}. Treating a tangential-field template as independent forcing
without a matching current repeats the realizability failure already flagged in
Chapter~1 \cite{stratton1941}.
Let \(\mathcal{M}_\omega\) denote the harmonic Maxwell operator of
Eq.~\eqref{eq:ch1.maxwell-operator-block}. The volume subsystem reads
\begin{equation}
\label{eq:ch2.maxwell-driven-operator}
\mathcal{M}_\omega[\mathbf{u}_\omega]=\mathbf{s}_\omega(\mathfrak{D}),
\end{equation}
and is satisfied when each residual row vanishes. Linearity in \(\mathbf{u}_\omega\)
follows from linearity of Eq.~\eqref{eq:ch1.maxwell-local}; spatially varying
\(\varepsilon(\mathbf{r})\) and \(\mu(\mathbf{r})\) modify coefficients inside
\(\mathcal{M}_\omega\) without changing the first-order block shape
(Section~\ref{sec:ch222}) \cite{stratton1941}. In uniform vacuum regions,
\(\mathcal{M}_\omega\) has translation-invariant symbol, the property exploited by
dyadic Green kernels in Section~\ref{sec:ch23}.
\begin{figure}[t]
\centering
\ChOneFigBlock{%
\begin{tikzpicture}[ch1fig, node distance=5mm and 7mm]
 \node[ch1box] (vol) {Volume\\$\mathcal{M}_\omega\mathbf{u}=\mathbf{s}$};
 \node[ch1box, below=5mm of vol] (bnd) {Boundary\\$\mathcal{B}_\Gamma\gamma_\Gamma\mathbf{u}=\mathbf{g}$};
 \node[ch1box, left=18mm of vol] (u) {State $\mathbf{u}_\omega$};
 \draw[ch1flow] (u.west) -- (u.west |- vol.east);
 \draw[ch1flow] (u.west |- vol.east) -- (vol.east);
 \draw[ch1flow] (u.east) -- (u.east |- bnd.west);
 \draw[ch1flow] (u.east |- bnd.west) -- (bnd.west);
\end{tikzpicture}%
}
\caption{Parallel volume and boundary channels for the closed harmonic problem:
\(\mathbf{u}_\omega\) simultaneously satisfies Maxwell residuals and prescribed
traces.}
\label{fig:ch2.maxwell-boundary-stack}
\end{figure}
~\eqref{eq:ch2.maxwell-driven-operator} and
\eqref{eq:ch1.boundary-data-operator}. A solver that enforces only volume rows
produces a formal solution compatible with Maxwell \emph{sources} but not necessarily
with termination physics; conversely, boundary data without volume balance violate
conservation. Subsection~\ref{sec:ch201} applies \(\mathcal{R}_\omega\) only after
both channels are closed \cite{harrington2001}. Tangential traces are extracted by the restriction
\begin{equation}
\label{eq:ch2.trace-restriction}
\gamma_\Gamma:\mathcal{F}_\Omega\to\mathcal{T}_\Gamma,
\qquad
\gamma_\Gamma[\mathbf{u}_\omega]=
\begin{bmatrix}
\left.\hat{\mathbf{n}}\times\E\right|_{\Gamma}\\
\left.\hat{\mathbf{n}}\times\H\right|_{\Gamma}
\end{bmatrix},
\end{equation}
with outward unit normal \(\hat{\mathbf{n}}\) on \(\Gamma\) and trace space
\(\mathcal{T}_\Gamma\) induced by interior states. The map
\(\gamma_\Gamma\) is not defined on arbitrary \(L^{2}\) six-tuples; Sobolev and
distribution refinements in Section~\ref{sec:ch211} specify its domain
\cite{stratton1941}. The closed operator stacks volume and boundary residuals:
\begin{equation}
\label{eq:ch2.maxwell-boundary-composite}
\mathcal{L}_\omega[\mathbf{u}_\omega]=
\begin{bmatrix}
\mathcal{M}_\omega[\mathbf{u}_\omega]\\
\mathcal{B}_\Gamma[\gamma_\Gamma[\mathbf{u}_\omega]]
\end{bmatrix}.
\end{equation}
A physically closed harmonic problem requires both volume Maxwell balance and boundary
termination to hold simultaneously; neither channel alone suffices
\cite{stratton1941}. The composite operator \(\mathcal{L}_\omega\) encodes that joint
requirement: its first block tests whether \((\E,\H)\) satisfies the curl--curl structure
with impressed sources removed to the forcing vector, while its second block tests whether
tangential traces on \(\Gamma\) match the declared termination (PEC, impedance, or outgoing)
\cite{harrington2001}. Experimentally, a solver that minimizes volume residuals while
ignoring port or chamber walls produces a formal Maxwell solution that may disagree with
what a network analyzer reports on the same structure \cite{stratton1941}. The harmonic
boundary-value problem is to find \(\mathbf{u}_\omega\in\mathcal{F}_\Omega\)
such that
\begin{equation}
\label{eq:ch2.closed-maxwell-bvp}
\mathcal{L}_\omega[\mathbf{u}_\omega]=
\begin{bmatrix}\mathbf{s}_\omega(\mathfrak{D})\\ \mathbf{g}_\Gamma\end{bmatrix}.
\end{equation}
Perfect conductors impose \(\hat{\mathbf{n}}\times\E=\mathbf{0}\); radiation and
Leontovich completions encode outgoing or dissipative relations on
\(\gamma_\Gamma[\mathbf{u}_\omega]\). The operator \(\mathcal{L}_\omega\) is the
Maxwell balance operator for this coupled system \cite{harrington2001}.
Hence $\mathcal{L}_\omega$ enforces Maxwell balance with impressed drive and boundary termination; without outgoing selection it does not distinguish export from storage, and every near-field probe reading $\mathcal{S}_\omega[\Jimp]$ presupposes its inversion on the declared domain \cite{stratton1941}.
Equation~\eqref{eq:ch2.maxwell-solution-map} is the resolvent relation
\begin{equation}
\label{eq:ch2.solution-map-resolvent}
\mathcal{S}_\omega[\Jimp]=\mathbf{u}_\omega
\quad\text{with}\quad
\mathcal{L}_\omega[\mathbf{u}_\omega]=
\begin{bmatrix}\mathbf{s}_\omega(\Jimp,\rho)\\ \mathbf{g}_\Gamma\end{bmatrix},
\end{equation}
where \(\rho\) is tied to \(\Jimp\) through
Eq.~\eqref{eq:ch1.charge-continuity-harmonic} when only electric impressed currents
drive the system. Invertibility of \(\mathcal{L}_\omega\) is not universal: existence
and uniqueness depend on \(\mathcal{B}_\Gamma\) and on whether outgoing selection is
imposed (Section~\ref{sec:ch243} for exterior domains; Section~\ref{sec:ch261} for
cavities) \cite{stratton1941}. \(\mathcal{S}_\omega\) is therefore a \emph{physical
branch} of \(\mathcal{L}_\omega^{-1}\), not an abstract algebraic inverse on all of
\(\mathcal{F}_\Omega\). \begin{figure}[t]
\centering
\ChOneFigBlock{%
\begin{tikzpicture}[ch1fig, node distance=10mm]
 \node[ch1box] (E) at (0,6mm) {$\curl\E$};
 \node[ch1box] (H) at (0,-6mm) {$\curl\H$};
 \node[ch1box, right=20mm of E] (r1) {$-\jj\omega\mu_0\H$};
 \node[ch1box, right=20mm of H] (r2) {$\Jimp+\jj\omega\varepsilon_0\E$};
 \draw[ch1flow] (E.east) -- (E.east |- r1.west);
 \draw[ch1flow] (E.east |- r1.west) -- (r1.west);
 \draw[ch1flow] (H.east) -- (H.east |- r2.west);
 \draw[ch1flow] (H.east |- r2.west) -- (r2.west);
 \draw[ch1flow, dashed] (r1.south) to[bend left=25] (H.east);
 \draw[ch1flow, dashed] (r2.north) to[bend left=25] (E.east);
\end{tikzpicture}%
}
\caption{Circulation coupling in \(\mathcal{M}_\omega\): each curl row feeds the
partner field; \(\Jimp\) appears as external drive on the magnetic row.}
\label{fig:ch2.first-order-curl-coupling}
\end{figure}
Single-field Helmholtz reductions equivalent to
Eqs.~\eqref{eq:ch1.vacuum-wave-source-driven}--\eqref{eq:ch1.vacuum-curl-second-order}
are used later when a scalar symbol simplifies kernels; the elimination belongs to
\ref{ch:01} and is not repeated. The electric-wave symbol used in
Chapter~\ref{ch:02} is recorded as
\begin{equation}
\label{eq:ch2.electric-wave-operator}
\mathcal{W}_{E,\omega}\E\equiv
\curl\!\left(\mu_0^{-1}\curl\E\right)-\omega^2\varepsilon_0\E
= -\jj\omega\mu_0\Jimp+\mu_0^{-1}\curl\mathbf{J}^{\mathrm{(mag)}}_{\mathrm{eff}},
\end{equation}
with \(\mathbf{J}^{\mathrm{(mag)}}_{\mathrm{eff}}\) collecting equivalent magnetic
drive when present \cite{stratton1941}. Subsection~\ref{sec:ch222} supplies the
magnetic counterpart and material-weighted generalizations. Variational structure is inherited from Section~\ref{sec:ch01.1.4}. For test state
\(\mathbf{v}_\omega=(\E',\H')\),
\begin{equation}
\label{eq:ch2.maxwell-energy-form}
\mathfrak{a}_\omega(\mathbf{u}_\omega,\mathbf{v}_\omega)=
\int_\Omega\Big[\mu_0\,\H\!\cdot\!\curl\E'-\varepsilon_0\,\E\!\cdot\!\curl\H'\Big]d^{3}r.
\end{equation}
Integration by parts connects Eq.~\eqref{eq:ch2.maxwell-energy-form} to surface
fluxes and to the weak formulation
\begin{equation}
\label{eq:ch2.maxwell-weak-form}
\mathfrak{a}_\omega(\mathbf{u}_\omega,\mathbf{v}_\omega)=\mathfrak{f}_\omega(\mathbf{v}_\omega)
\quad\forall\,\mathbf{v}_\omega\in\mathcal{F}_\Omega^{0},
\end{equation}
where \(\mathcal{F}_\Omega^{0}\) carries homogeneous boundary traces compatible with
\(\mathcal{B}_\Gamma\) and \(\mathfrak{f}_\omega\) is induced by
\(\mathbf{s}_\omega(\mathfrak{D})\) \cite{poynting1884}. Equation
\eqref{eq:ch2.maxwell-weak-form} is the variational face of
Eq.~\eqref{eq:ch2.closed-maxwell-bvp}; finite-element discretizations favor this
route when material coefficients are piecewise constant
\cite{harrington2001}.
Well-posedness of Eq.~\eqref{eq:ch2.maxwell-weak-form} is not automatic on
unbounded domains or on multiply connected cavities. On a bounded region with
passive \(\mathcal{B}_\Gamma\) and homogeneous traces in
\(\mathcal{F}_\Omega^{0}\), coercivity fails only along divergence-free
cavity modes that satisfy both curl-curl balance and homogeneous tangential
data---precisely the null space that Section~\ref{sec:ch214} isolates before
imposing gauge or polarization fixes \cite{harrington2001}. For
exterior problems, coercivity is recovered only after restricting to the
outgoing subspace selected by Sommerfeld-type completion; without that restriction,
the bilinear form is indefinite and the weak problem admits non-physical
growing solutions. The practical rule for antenna analysis is therefore
threefold: declare \(\mathcal{D}(\mathcal{L}_\omega)\), specify
\(\mathcal{B}_\Gamma\) on every connected piece of \(\Gamma\), and confirm that
the impressed tuple lies in the compatible range of
Eq.~\eqref{eq:ch2.source-compatibility}. Numerical codes that enforce volume
Maxwell rows on truncated domains but leave exterior faces open implement an
\emph{implicit} \(\mathcal{B}_\Gamma\); unless that implicit completion matches
the laboratory termination, \(\mathcal{S}_\omega\) in the code is a different
operator from the one measured on the bench \cite{harrington2001}. Domain specification cannot remain at \(C^\infty\) when sources are distributional.
The minimal energy class declared here is
\begin{equation}
\label{eq:ch2.operator-domain-declaration}
\mathcal{D}(\mathcal{L}_\omega)\subset
\big\{\mathbf{u}_\omega=(\E,\H):
\E,\H\in L^{2}(\Omega)^{3},\ \curl\E,\curl\H\in L^{2}(\Omega)^{3}\big\},
\end{equation}
where $\mathcal{L}_\omega$ is the composite Maxwell operator on the six-component state
$\mathbf{u}_\omega=(\E,\H)$, $\mathcal{D}(\mathcal{L}_\omega)$ its domain of square-integrable
curl fields compatible with $\mathcal{B}_\Gamma$, and $\Omega$ the analysis region
\cite{stratton1941}. Section~\ref{sec:ch211} refines embeddings and compactness; evaluating
Eq.~\eqref{eq:ch2.maxwell-driven-operator} outside
\(\mathcal{D}(\mathcal{L}_\omega)\) generates spurious boundary terms that corrupt
flux identities \cite{stratton1941}.
Piecewise smooth \(\varepsilon(\mathbf{r})\) and \(\mu(\mathbf{r})\) modify
coefficients inside \(\mathcal{M}_\omega\) while preserving the block template of
Eq.~\eqref{eq:ch1.maxwell-operator-block}. Interface jumps satisfy
Eq.~\eqref{eq:ch1.interface-conditions}; multi-region Green operators in
Section~\ref{sec:ch234} realize \(\mathcal{L}_\omega^{-1}\) on composite domains
\cite{stratton1941}. The coupling structure described above extends naturally to practical
terminal conditions. Finite feed structures and external circuits frequently impose
impedance relations between tangential field components on the boundary
\cite{harrington2001}:
\begin{equation}
\label{eq:ch2.boundary-impedance-operator}
\mathcal{Z}_\Gamma\!
\begin{bmatrix}
\left.\hat{\mathbf{n}}\times\E\right|_{\Gamma}\\
\left.\hat{\mathbf{n}}\times\H\right|_{\Gamma}
\end{bmatrix}=\mathbf{0},
\end{equation}
In Eq.~\eqref{eq:ch2.boundary-impedance-operator},
\(\mathcal{Z}_\Gamma\) is a linear boundary impedance operator (SI units absorbed in its
definition) acting on the tangential Maxwell traces
\(\hat{\mathbf{n}}\times\E|_{\Gamma}\) and \(\hat{\mathbf{n}}\times\H|_{\Gamma}\) on
\(\Gamma=\partial\Omega\); admissible states satisfy the terminal or absorbing relation
\(\mathcal{Z}_\Gamma[\text{traces}]=\mathbf{0}\).
encoding Leontovich, absorbing, or mode-matched completions \cite{harrington2001}.
When Eq.~\eqref{eq:ch2.boundary-impedance-operator} replaces explicit
\(\mathbf{g}_\Gamma\) in Eq.~\eqref{eq:ch2.closed-maxwell-bvp}, admissible states
satisfy the impedance manifold on \(\Gamma\). Exterior outgoing problems correspond
to a limit in which \(\mathcal{Z}_\Gamma\) selects the Sommerfeld branch of
Section~\ref{sec:ch243}; conflating a finite-chamber termination with that limit
without normalization corrupts far-zone patterns \cite{stratton1941}.
Forcing must lie in the range implied by charge conservation. Divergence of the
Amp\`ere row together with Eq.~\eqref{eq:ch1.charge-continuity-harmonic} yields
\begin{equation}
\label{eq:ch2.source-compatibility}
\divg\mathbf{s}_{\omega,2}+\jj\omega\,\mathbf{s}_{\omega,3}=0
\end{equation}
on the second and third blocks of Eq.~\eqref{eq:ch2.maxwell-forcing-vector}.
Inconsistent \((\Jimp,\rho)\) pairs lie outside physically realizable drive;
numerical solvers may still return vectors that violate
Eq.~\eqref{eq:ch2.source-compatibility} \cite{stratton1941}. A compact
current filament with vanishing \(\rho\) automatically satisfies
Eq.~\eqref{eq:ch2.source-compatibility} because \(\mathbf{s}_{\omega,3}=0\).
When a volumetric \(\rho\) is prescribed independently of \(\Jimp\), the pair
must be adjusted so that their divergence balance matches the Amp\`ere--Gauss
rows; otherwise the closed problem is mathematically posed but physically
inconsistent. Feed models that inject \(\Jimp\) without updating \(\rho\) are
common in method-of-moments implementations; they remain valid only when the
discretization enforces continuity on the impressed grid or when the problem
is magnetostatic in the relevant band \cite{harrington2001}.
Linking to Subsection~\ref{sec:ch201}, radiation acts on closed solutions:
\begin{equation}
\label{eq:ch2.maxwell-radiation-composition}
\mathcal{R}_\omega=\Pi_{\mathrm{rad}}\circ\mathcal{S}_\omega
=\Pi_{\mathrm{rad}}\circ\mathcal{L}_\omega^{-1}\circ
\begin{bmatrix}\mathbf{s}_\omega(\cdot)\\ \mathbf{g}_\Gamma\end{bmatrix},
\end{equation}
with \(\mathcal{L}_\omega^{-1}\) understood on the outgoing-admissible branch fixed by
\(\mathcal{B}_\Gamma\) and with \(\Pi_{\mathrm{rad}}\) from
Eq.~\eqref{eq:ch2.radiation-operator-def} \cite{stratton1941}. Completion must precede
radiative selection: applying \(\Pi_{\mathrm{rad}}\) to non-closed states mixes
reactive content into far-zone observables \cite{harrington2001}.
The energy form supports an integration-by-parts identity
\begin{equation}
\label{eq:ch2.maxwell-adjoint-pairing}
\mathfrak{a}_\omega(\mathbf{u}_\omega,\mathbf{v}_\omega)=
\langle\mathcal{M}_\omega[\mathbf{u}_\omega],\mathbf{v}_\omega\rangle_{\Omega}
+\int_\Gamma\big(\hat{\mathbf{n}}\times\E\cdot\H'-\hat{\mathbf{n}}\times\E'\cdot\H\big)dS,
\end{equation}
where \(\langle\cdot,\cdot\rangle_{\Omega}\) is the componentwise \(L^{2}\) pairing
\cite{stratton1941}. Self-adjoint truncations on bounded domains require that
\(\mathcal{B}_\Gamma\) cancel the boundary integrand on admissible traces; criteria
are supplied in Section~\ref{sec:ch214}. Equation
\eqref{eq:ch2.maxwell-adjoint-pairing} is the local spine for those arguments.
Fredholm structure note: for bounded \(\Omega\) with passive \(\mathcal{B}_\Gamma\), \(\mathcal{L}_\omega\) is a Fredholm operator on
\(\mathcal{D}(\mathcal{L}_\omega)\) modulo cavity modes, with index tied to
topology and boundary class. Exterior domains replace discrete spectra with
continuous radiation modes, and \(\mathcal{L}_\omega^{-1}\) becomes a resolvent on
the outgoing subspace rather than an ordinary inverse. Section~\ref{sec:ch263}
develops the spectral/Fredholm foundation; here the point is that invertibility questions
for \(\mathcal{S}_\omega\) are questions about \(\mathcal{L}_\omega\) on the
\emph{declared} domain, not about formal elimination of \(\H\) alone
\cite{harrington2001}.
Linearity of \(\mathcal{L}_\omega\) in \(\mathbf{u}_\omega\) implies that
superposition holds for fixed boundary class: if
\(\mathcal{L}_\omega[\mathbf{u}_\omega^{(a)}]=[\mathbf{s}^{(a)},\mathbf{g}^{(a)}]\)
and \(\mathcal{L}_\omega[\mathbf{u}_\omega^{(b)}]=[\mathbf{s}^{(b)},\mathbf{g}^{(b)}]\)
with identical \(\mathcal{B}_\Gamma\), then
\(\mathcal{L}_\omega[\mathbf{u}_\omega^{(a)}+\mathbf{u}_\omega^{(b)}]=
[\mathbf{s}^{(a)}+\mathbf{s}^{(b)},\mathbf{g}^{(a)}+\mathbf{g}^{(b)}]\).
Changing \(\mathcal{B}_\Gamma\) while holding \(\Jimp\) fixed breaks this
principle because boundary rows are part of the operator, not of the drive.
Laboratory comparisons that swap terminations without re-solving the coupled
system therefore compare different \(\mathcal{L}_\omega\), not different
``modes'' of the same map \cite{stratton1941}. The radiation
operator \(\mathcal{R}_\omega\) inherits the same sensitivity: composition
Eq.~\eqref{eq:ch2.maxwell-radiation-composition} is linear in impressed data
only after \(\mathcal{B}_\Gamma\) and the outgoing branch are held fixed
\cite{stratton1941}.
Subsection~\ref{sec:ch202} has packaged harmonic Maxwell theory as the composite operator
\(\mathcal{L}_\omega\), recorded its resolvent relation to \(\mathcal{S}_\omega\), and
linked radiation composition Eq.~\eqref{eq:ch2.maxwell-radiation-composition} to operator
inversion on the outgoing-admissible branch. Subsection~\ref{sec:ch203} records how those
maps inherit the admissibility, regime, and measurement criteria established in
\ref{ch:01}.
